\section{Method}\label{sec:method}

\subsection{Thermal model}
Let $\vec{x} \in \R^2$ be a ground position and \(\heat(\vec{x})\) the measured surface temperature. A sheep appears as a warm blob; we write $\pair{u}{v}$ for a pixel and $\norm{\vec{x} - \vec{y}}$ for ground distance. The background is modelled as
\[
  \heat_0(\vec{x}) = a + \inner{\vec{b}}{\vec{x}}, \qquad a \in \R,\ \vec{b} \in \R^2 ,
\]
which is enough on the short flights we use. An unnumbered display also appears in dollars: $$ r = \unit[px]{12} $$ and the text continues after it.

\begin{equation}
  \lambda(\vec{x}) = \sum_{i=1}^{N} k\bigl(\norm{\vec{x} - \vec{s}_i}\bigr)
  \label{eq:intensity}
\end{equation}
where $k$ is a smooth kernel. The count estimate is
\begin{equation*}
  \hat{N} = \argmin_{N \ge 0} \; \E\left[ \bigl( \census(\heat) - N \bigr)^2 \right] .
\end{equation*}

\begin{align}
  L(\theta) &= \sum_{t=1}^{T} \log p(\heat_t \mid \theta) \label{eq:lik} \\
            &= -\frac{1}{2\sigma^2} \sum_{t=1}^{T} \norm{\heat_t - \mu_\theta}^2 + C , \nonumber \\
  \hat\theta &= \argmin_\theta \, -L(\theta) . \tag{MLE}
\end{align}

\begin{align*}
  a &= b + c \\
  d &= e
\end{align*}

\begin{gather}
  x^2 + y^2 = r^2 \\
  e^{i\pi} + 1 = 0
\end{gather}

\begin{multline}
  f(a, b, c) = a^3 + b^3 + c^3 - 3abc \\
  = (a + b + c)(a^2 + b^2 + c^2 - ab - bc - ca)
\end{multline}

\begin{equation}
\begin{split}
  S &= \sum_{i} w_i \\
    &= 1
\end{split}
\end{equation}

\begin{subequations}\label{eq:pair}
\begin{equation}
  \dot{p} = -\alpha p + \beta q
\end{equation}
\begin{equation}
  \dot{q} = \alpha p - \beta q
\end{equation}
\end{subequations}

The disturbance indicator is piecewise,
\begin{equation}
  D_t = \begin{cases}
    1 & \text{if more than 5\% of animals stand at time } t, \\
    0 & \text{otherwise,}
  \end{cases}
  \qquad
  M = \begin{pmatrix} 1 & 0 \\ 0 & \gamma \end{pmatrix},
  \quad
  B = \begin{bmatrix} a & b \\ c & d \end{bmatrix} .
\end{equation}

\begin{eqnarray}
  y &=& m x + c \\
  z &=& y^2
\end{eqnarray}

An equation that KaTeX cannot typeset must degrade gracefully:
\begin{equation}
  \sideset{_a^b}{'}{\sum}_{i \in \mathcal{H}} w_i \heat_i
\end{equation}

\subsection{Guarantees}

\begin{definition}[Huddle]\label{def:huddle}
A \term{huddle} is a maximal set of animals whose thermal outlines are connected at the threshold $\wake$.
\end{definition}

\begin{theorem}[Consistency]
\label{thm:consistency}
Suppose the kernel $k$ is bounded and the flight height is fixed. Then $\hat{N} \to N$ in probability as the number of passes grows.
\end{theorem}

\begin{lemma}
Huddles of size at most three are separable.
\end{lemma}

\begin{proof}
By \cref{def:huddle}, a huddle of size $m \le 3$ has at most $m$ local maxima of $\lambda$ in \eqref{eq:intensity}. Counting maxima gives
\[
  \census(\heat) = m ,
\]
which proves the claim.
\end{proof}

\begin{proposition}
The disturbance score is monotone in flight height.
\end{proposition}

\begin{remark}
The converse of \cref{thm:consistency} is false; see \cref{app:proofs}.
\end{remark}

\begin{flockrule}
Never fly upwind of a sleeping flock.
\end{flockrule}

\begin{quote}
``You can fool a ewe once.'' --- a shepherd at Dunmore, on the first night of flights.
\end{quote}

\begin{quotation}
A longer quotation runs to several sentences. It is set as its own paragraph and is editable like any other prose.

It even has a second paragraph.
\end{quotation}
